%!TEX root = ../chapter07.tex 
%----------------------------------------------------------------------------------------
%	
%----------------------------------------------------------------------------------------
%----------------------------------------------------------------------------------------
%	 
%----------------------------------------------------------------------------------------


%----------------------------------------------------------------------------------------
%	EXERCICES
%----------------------------------------------------------------------------------------

\newpage 
\loadgeometry{widelayout}  
\pagestyle{scrheadings} % Print headers again  
\KOMAoptions{
	headwidth=textwithmarginpar,
	headsepline=true,
	headsepline= 0.5pt,
	footwidth=textwithmarginpar,
} 
\onehalfspacing
\section{Exercices}
\begin{exercise}[Vérifier si une valeur est solution d'une équation à 1 inconnue]\hfill 
	
\noindent\QCM[VF,Alterne,Stretch=1.0, Largeur=0.75cm]{%
	{$70$ est solution de l'équation $x^{300}=1$ d'inconnue $x$}&2,%
	{$2$ est une solution de l'équation $2x+1=5$ d'inconnue $x$}&1,%
	{$2$ est la solution de l'équation $2x+1=5$ d'inconnue $x$}&1,%
	{$3$ est une solution de l'équation $x^2=3x$ d'inconnue $x$}&1,%
	{$3$ est la solution de l'équation $x^2=3x$ d'inconnue $x$}&2,%
	{$3$ est une solution de l'équation $\abs{x-5}=2$ inconnue $x$}&1,%
	{$-3$ est une solution de l'équation $\abs{x-5}=2$ inconnue $x$}&2,%
	{$-\sqrt{3}$ est une solution de l'équation $3x^2=-9$ inconnue $x$}&2,%
	{$-\sqrt{3}$ est une solution de l'équation $x^2+2=5$ inconnue $x$}&1,%
%	{$\sqrt{3}$ est une solution de l'équation $(3x-1)(x-\sqrt{3})=0$ inconnue $x$}&1,%
	{$\sqrt{3}$ est une solution de l'équation $(2-x)(2+x)=1$ inconnue $x$}&1,% 
} 
\end{exercise}
\begin{exercise}[Vérifier si un couple est solution d'une équation à 2 inconnues]\hfill 
	
\noindent\QCM[VF,Alterne,Stretch=1.0, Largeur=0.75cm]{%
	{$x=-3$  et $y=3$ est couple solution de l'équation $6x+3y=-2$ d'inconnue  $(x;y)$}&2,%
	{$x=-\tfrac{2}{3}$  et $y=\tfrac{2}{3}$ est couple solution de l'équation $6x+3y=-2$ d'inconnue  $(x;y)$}&1,% 
	{$x=0$  et $y=-1$ est couple solution de l'équation $-9x-7y=-7$ d'inconnue  $(x;y)$}&2,% 
	{$x=\tfrac{7}{9}$  et $y=0$ est couple solution de l'équation $-9x-7y=-7$ d'inconnue  $(x;y)$}&1,% 
	{$x=18$  et $y=\tfrac{323}{20}$ est couple solution de $17x-20y=-17$ d'inconnue  $(x;y)$}&1,% 
	{$x=2$  et $y=5$ est un couple solution de $(x-2)(y+5)=0$ d'inconnue  $(x;y)$}&1,%
	{$x=3$  et $y=-5$ est un couple solution de  $(x-2)(y+5)=0$ d'inconnue  $(x;y)$}&1,% 
	{$x=3$  et $y=\tfrac{1}{3}$ est le couple solution de l'équation $xy=1$ d'inconnue  $(x;y)$}&2,% 
	{$x=0$  et $y=0$ est le couple solution de l'équation $xy=0$ d'inconnue  $(x;y)$}&2,% 
	{$x=-6$  et $y=8$ est un couple solution de $x^2+y^2=10^2$ d'inconnue  $(x;y)$}&1,% 
} 
\end{exercise}
\begin{exercise} \hfill 
	
On considère la résolution suivante :  $\begin{WithArrows}
	2(x+6)-4&=8+6x \Arrow{On développe le membre de gauche}\\
	2x+12-4&=8+6x\Arrow{On simplifie le membre de gauche}\\
	2x+8&=8+6x\Arrow{On ajoute $-8$ aux 2 membres}\\
	2x&=6x\Arrow{On divise les 2 membres par $x$}\\
	2&=6
\end{WithArrows}$ 
 
Quelle réponse explique l'erreur à l'étape \og on divise les 2 membres par $x$ \fg{} :
\begin{enumerate}[leftmargin=*,label={(\Alph*)}]
	\item Diviser par $x$ les 2 membres de $2x=6x$ donne $2x^2=6x^2$.
	\item Diviser par $x$ les 2 membres de $2x=6x$ donne $2=6x$.
	\item Diviser par $x$ les 2 membres de $2x=6x$ laisse croire que l'équation n'a pas de solutions alors que $x=0$ est solution.
	\item Diviser par $x$ les 2 membres de $2x=6x$ laisse croire que l'équation n'a pas de solutions alors que $x=3$ est solution.
\end{enumerate}  
\end{exercise}
\begin{exercise}\hfill 
	
	 Compléter la phrase pour expliquer l'erreur dans la résolution suivante : \hfill 
	
\noindent\parbox{1\linewidth}{
$\begin{WithArrows}
	x(x-1)-3(x-1)  & =5x-5  \Arrow{On factorise les deux membres, facteurs commun $(x-1)$}\\
	(x-3) (x-1)  & =5(x-1)  \Arrow{On divise les deux membres par $\divisionsymbol (x-1)$}\\
	(x-3) & = 5  \Arrow{On ajoute aux deux membres $3$}\\
	x & =8  
\end{WithArrows}$ 
}

Diviser par $x-1$ les 2 membres de $(x-3)(x-1)=5(x-1)$ laisse croire que l'équation a une solution unique $x=\ldots$ , alors que $x= \ldots $ est aussi solution de l'équation de départ ! 
\end{exercise} 
\begin{tBox}
	On change pas les solutions d'une équation si :
	\begin{enumerate}[label=\textbullet, leftmargin=*]
		\item on \textbf{additionne, soustrait} aux \textbf{2 membres} de l'équation  une même expression.
		\item on \textbf{multiplie/divise} les \textbf{2 membres} de l'équation par \textbf{une même}  expression \textbf{non nulle}.
		\item on \textbf{développe, factorise, réduit ...} un des deux membres de l'équation.
	\end{enumerate}  
\end{tBox} 


\begin{example}  Résoudre dans $\mathbb{R}$ les équations suivantes :\hfill 
	
\noindent \parbox{0.45\linewidth}{	
	$\begin{WithArrows}
		4(x+1)&=x+7 \\
		&= \rule{0pt}{1.3em} \\
		& = \rule{0pt}{1.3em}  \\
		& = \rule{0pt}{1.3em}  \\
		& = \rule{0pt}{1.3em}  
	\end{WithArrows}$ 	
}\hfill \parbox{0.45\linewidth}{	
	$\begin{WithArrows}
		4(x-1)&=-7x+5 \\
		&= \rule{0pt}{1.3em} \\
		& = \rule{0pt}{1.3em}  \\
		& = \rule{0pt}{1.3em}  \\
		& = \rule{0pt}{1.3em} 
	\end{WithArrows}$ 
}  
\end{example} 
\begin{exercise}\label{chapter07:exercice:equationA} Résoudre dans $\mathbb{R}$ les équations suivantes d'inconnue $x$: \hfill 
	\begin{multicols}{3}
 		\begin{enumerate}[label={$(E_{\arabic*})$}, leftmargin=*]
			\item $3(x+5)=x+1$
			\item $3(x+5)=3(2x+1)$
			\item $3(x+5)-3(2x-1)=0$
			\item $3(2x-5)=3(x-1)$
			\item $3(x-5)=-3(2x+1)$ 
			\item $-3(2x+5)=-3(2x+1)$
			\item $-2(3x-6)=3(-2x+4)$
			\item $-3(x-5)=4-3(2x-1)$
			\item $-3(2x-1)=x-3(x-5)$ 
			%		\item[$E_{1}\colon $] $x+5=3(x+1)$  
			%		\item[$E_{5}\colon$]  $3(x-5)=3(2x-1)$
			%		\item[$E_{7}\colon$]  $ 3(x-5)=3(2x+1)$
		\end{enumerate}
	\end{multicols} 	
\end{exercise}

\begin{example}  Résoudre dans $\mathbb{R}$ les équations suivantes d'inconnue $x$: 
	
\noindent \parbox{0.45\linewidth}{	
	$\begin{WithArrows}
		\tfrac{x+2}{2}+\tfrac{x-4}{5}&=6 \\
		&= \rule{0pt}{1.3em} \\
		& = \rule{0pt}{1.3em}  \\
		& = \rule{0pt}{1.3em}  \\
		& = \rule{0pt}{1.3em}  
	\end{WithArrows}$ 	
}\hfill \parbox{0.45\linewidth}{	
	$\begin{WithArrows}
		\tfrac{x+2}{2}-\tfrac{x-4}{6}&=6 \\
		&= \rule{0pt}{1.3em} \\
		& = \rule{0pt}{1.3em}  \\
		& = \rule{0pt}{1.3em}  \\
		& = \rule{0pt}{1.3em} 
	\end{WithArrows}$ 
}   
\end{example}
%\newpage
\begin{exercise}\label{chapter07:exercice:equationB} Résoudre dans $\mathbb{R}$ les équations suivantes d'inconnue $x$: \hfill  
	\begin{multicols}{3}
 		\begin{enumerate}[label={$(E_{\arabic*})$}, leftmargin=*]
			\item $	\tfrac{x+2}{2}+\tfrac{x+4}{3}=4$
			\item $	\tfrac{x+2}{2}+\tfrac{x+4}{8}=4x$
			\item $	\tfrac{x+2}{6}+\tfrac{x+4}{8}=9$
			\item $ \tfrac{x+2}{6}-\tfrac{x+4}{8}=9$
			\item $	\tfrac{x-2}{6}-\tfrac{x-4}{8}=9$
			\item $\tfrac{2x-1}{3}-\tfrac{x}{2}=3x$
			\item $	\tfrac{4x-4}{8}-\tfrac{4x-2}{6}=9$
			\item $\tfrac{x}{2}-\tfrac{2x-1}{3}=1$
			\item $\tfrac{x-5}{3}-2(x-7)=2x+1$
			%			\item[$E_{3}\colon$]  $	\frac{x+2}{2}+\frac{x+4}{8}=9$
			%			\item[$E_{6}\colon$] $ \frac{x+2}{6}-\frac{x-4}{8}=9$ 
			%			\item[$E_{8}\colon$] $\frac{2x-2}{2}-\frac{2x-4}{8}=9$
			%			\item[$E_{11}\colon$] $	\frac{4x-4}{8}-\frac{4x-2}{6}=1$
		\end{enumerate}
	\end{multicols} 	
\end{exercise}
%----------------------------------------------------------------------------------------
%	 
%----------------------------------------------------------------------------------------
\newpage 
%\section{Exercices isoler l'inconnue pour résoudre}
\begin{example} Résoudre dans $\mathbb{R}$ les équations suivantes d'inconnue $x$.
	
	\noindent \parbox{0.25\linewidth}{	
		$\begin{WithArrows}
			\abs{x}&=  10 \\ 
			& \rule{0pt}{1.3em}  \\
			&  \rule{0pt}{1.3em} \\
			&   \rule{0pt}{1.3em}  
		\end{WithArrows}$ 	
	}\hfill \parbox{0.25\linewidth}{	
		$\begin{WithArrows}
			\abs{2x-5}&=  10 \\ 
			& \rule{0pt}{1.3em}  \\
			&  \rule{0pt}{1.3em} \\
			&   \rule{0pt}{1.3em}  
		\end{WithArrows}$ 	
	}\hfill \parbox{0.40\linewidth}{	
		$\begin{WithArrows}
			\abs{2x-7}&=\abs{x+3} \\ 
			& \rule{0pt}{1.3em}  \\
			&  \rule{0pt}{1.3em} \\
			&   \rule{0pt}{1.3em}  
		\end{WithArrows}$ 	
	} 
\end{example}
\begin{exercise}\label{chapter07:exercice:equationC} Résoudre dans $\mathbb{R}$ les équations suivantes d'inconnue $x$: \hfill  
	\begin{multicols}{3}
		\begin{enumerate}[label={$(E_{\arabic*})$}, leftmargin=*]
			\item  $\abs{2x}=3$
			\item  $\abs{3x-5}=7$
			\item  $\abs{4x-3}=-1$
			\item  $\abs{5-2x}=10^{-2}$
			\item  $\abs{2x-1}=\abs{2x-7}$
			\item  $\abs{5x+1}=\abs{2x+5}$
		\end{enumerate}
	\end{multicols} 	
\end{exercise}
\begin{example} Résoudre dans $\mathbb{R}$ les équations suivantes d'inconnue $x$.
	
	\noindent \parbox{0.22\linewidth}{	
		$\begin{WithArrows}
			3\abs{x} + 5 &=  14 \\ 
			& \rule{0pt}{1.3em}  \\
			&  \rule{0pt}{1.3em} \\
			&   \rule{0pt}{1.3em}  \\
			&   \rule{0pt}{1.3em}  
		\end{WithArrows}$ 	
	}\hfill \parbox{0.22\linewidth}{	
		$\begin{WithArrows}
			3x^2 + 5 &=  14 \\ 
			& \rule{0pt}{1.3em}  \\
			&  \rule{0pt}{1.3em} \\
			&   \rule{0pt}{1.3em}  \\
			&   \rule{0pt}{1.3em}  
		\end{WithArrows}$ 	
	} \hfill \parbox{0.22\linewidth}{	
		$\begin{WithArrows}
			3\abs{x} + 5 &=  4 \\ 
			& \rule{0pt}{1.3em}  \\
			&  \rule{0pt}{1.3em} \\
			&   \rule{0pt}{1.3em}  \\
			&   \rule{0pt}{1.3em}  
		\end{WithArrows}$ 	
	}\hfill \parbox{0.22\linewidth}{	
		$\begin{WithArrows}
			3x^2 + 5 &=  4 \\ 
			& \rule{0pt}{1.3em}  \\
			&  \rule{0pt}{1.3em} \\
			&   \rule{0pt}{1.3em}   \\
			&   \rule{0pt}{1.3em} 
		\end{WithArrows}$ 	
	} 
\end{example}
\begin{exercise}\label{chapter07:exercice:equationD} Résoudre dans $\mathbb{R}$ les équations suivantes d'inconnue $x$: \hfill  
	 
	\begin{multicols}{3}
		\begin{enumerate}[label={$(E_{\arabic*})$}, leftmargin=*]
			\item  $5\abs{x}-3=6$
			\item  $-2\abs{x}+11=6$
			\item  $5x^2-7=-1$
			\item  $-2x^2+5=-9$
			\item  $-3\abs{x}+2=6$
			\item  $x^2+5=0$
		\end{enumerate}
	\end{multicols} 	
\end{exercise}
\begin{example}[Isoler une variable]  
Pour chacune des relations suivantes, exprimer $x$ en fonction des autres variables. On supposera que les lettres prennent des valeurs strictement positives.

\noindent \parbox{0.16\linewidth}{	
	$\begin{WithArrows}
		A&=\tfrac{Bx}{2}  \\ 
		& \rule{0pt}{1.3em}  \\
		&  \rule{0pt}{1.3em} \\
		&   \rule{0pt}{1.3em} \\
		&   \rule{0pt}{1.3em} 
	\end{WithArrows}$ 	
}\hfill \parbox{0.16\linewidth}{	
	$\begin{WithArrows}
		P&=2\pi x -3\\ 
		& \rule{0pt}{1.3em}  \\
		&  \rule{0pt}{1.3em} \\
		&   \rule{0pt}{1.3em} \\
		&   \rule{0pt}{1.3em} 
	\end{WithArrows}$ 	
} \hfill \parbox{0.19\linewidth}{	
	$\begin{WithArrows}
		2y-1&=\frac{3}{x} \\ 
		& \rule{0pt}{1.3em}  \\
		&  \rule{0pt}{1.3em} \\
		&   \rule{0pt}{1.3em}  \\
		&   \rule{0pt}{1.3em}
	\end{WithArrows}$ 	
}\hfill \parbox{0.19\linewidth}{	
	$\begin{WithArrows}
		2y+3&=5x+3 \\ 
		& \rule{0pt}{1.3em}  \\
		&  \rule{0pt}{1.3em} \\
		&   \rule{0pt}{1.3em}  \\
		&   \rule{0pt}{1.3em}
	\end{WithArrows}$ 	
} \hfill \parbox{0.22\linewidth}{	
$\begin{WithArrows}
	yx+5&=2x-3 \\ 
	& \rule{0pt}{1.3em}  \\
	&  \rule{0pt}{1.3em} \\
	&   \rule{0pt}{1.3em} \\
	&   \rule{0pt}{1.3em}  
\end{WithArrows}$ 	
} 	 
\end{example}
\begin{exercise}[à vous] \label{chapter07:exercice:equationE} Mêmes consignes

\begin{multicols}{4}
	\begin{enumerate}[label={\alph*)\rule{0pt}{1.6em}},leftmargin =*] 
		\item $2\pi x  =A$  
		\item $\tfrac{2\pi y}{x}  =A$    
		\item $3x-5=2y$ 
		\item $5x+7y=1$
		\item $-3=2x-9y$
		\item $12=2y+5x$
		\item $5y+3=\tfrac{2}{x}$
		\item $7=3y-\tfrac{5}{x}$
		\item $7y=2xy+5$ 
		\item $2yx+3= 3x+7$
		\item $5yx+3= 7x+2y$
		\item $2yx-3= 3x-5y$
	\end{enumerate}
\end{multicols}
\end{exercise} 
%----------------------------------------------------------------------------------------
%	 
%----------------------------------------------------------------------------------------

\newpage
\begin{proof}[solution de l'exercice \ref{chapter07:exercice:equationA}]\hfill 
	\begin{multicols}{5}
		\begin{enumerate}[label={$(E_{\arabic*}$)}, leftmargin=*]
			\item $\mathcal{S}=\{-7\}$
			\item $\mathcal{S}=\{4\}$ 
			\item $\mathcal{S}=\{6\}$ 
			\item   $\mathcal{S}=\{4\}$ 
			\item $\mathcal{S}=\{\tfrac{12}{9}\}$  
			\item {\fontspec{Segoe UI Emoji} 😱}
			\item {\fontspec{Segoe UI Emoji} 🤔}
			\item $\mathcal{S}=\{-\tfrac{8}{3}\}$ 
			\item $\mathcal{S}=\{-3\}$ 
			%			\item[$E_{1}\colon $] $\mathcal{S}=\{1\}$  
			%			\item[$E_{5}\colon$] $\mathcal{S}=\{-4\}$ 
			%			\item[$E_{7}\colon$] $\mathcal{S}=\{ -6\}$ 
		\end{enumerate}
	\end{multicols} 	
\end{proof} 
\begin{proof}[solutions de l'exercice \ref{chapter07:exercice:equationB}]\hfill 
	\begin{multicols}{5}
		\begin{enumerate}[label={$(E_{\arabic*}$)}, leftmargin=*]
			\item  $\mathcal{S}=\{2\}$ 
			\item $\mathcal{S}=\{\tfrac{4}{9}\}$
			\item $\mathcal{S}=\{28\}$ 
			\item $\mathcal{S}=\{220\}$ 
			\item $\mathcal{S}=\{ 212\}$ 
			\item $\mathcal{S}=\{-\tfrac{2}{17}\}$  
			\item $\mathcal{S}=\{-55\}$ 
			\item $\mathcal{S}=\{-4\}$  
			\item $\mathcal{S}=\{\tfrac{34}{11}\}$
			%			\item[$E_{3}\colon$] $\mathcal{S}=\{12\}$ 
			%			\item[$E_{6}\colon$]  $\mathcal{S}=\{196\}$ 
			%			\item[$E_{8}\colon$] $\mathcal{S}=\{106 \}$   
			%			\item[$E_{11}\colon$]  $\mathcal{S}=\{-7\}$  
		\end{enumerate}
	\end{multicols} 	
\end{proof}

\begin{proof}[solutions de l'exercice \ref{chapter07:exercice:equationC}]\hfill  
\begin{pycode}
from re import sub
import sympy as sp
from sympy.core.sympify import kernS

x, y, z, t , a, b = sp.symbols('x y z t a b', real=True)
k, m, n = sp.symbols('k m n', integer=True)
f, g, h = sp.symbols('f g h', cls=sp.Function)

# Create a list of functions to include in the table
funcs = [
'abs(2*x)-3',
'abs(2*x-5)-7',
'abs(4*x-3)+1',
'abs(5-2*x)-10**(-2)',
'abs(2*x-1)-abs(2*x-7)',
'abs(5*x+1)-abs(2*x+5)',
]
n=1
print(r'\begin{multicols}{4} \begin{enumerate}[label={$S_{\arabic*}=$}] ')
for func in funcs:
	print('\item ')
	solution = set(sp.solve(eval(func), x))
	string =  '$'+   sp.latex(solution ) + '$; '
	print(string)
print(r'\end{enumerate} \end{multicols} ')
\end{pycode}
\end{proof}

\begin{proof}[solutions de l'exercice \ref{chapter07:exercice:equationD}]\hfill  
\begin{pycode}
from re import sub
import sympy as sp
from sympy.core.sympify import kernS

x, y, z, t , a, b = sp.symbols('x y z t a b', real=True)
k, m, n = sp.symbols('k m n', integer=True)
f, g, h = sp.symbols('f g h', cls=sp.Function)

# Create a list of functions to include in the table
funcs = [
'5*abs(x)-9', 
'-2*abs(x)+5',
'5*x**2-6',
'-2*x**2+14',
'-3*abs(x)-4',
'x**2+5', 
]
n=1
print(r'\begin{multicols}{4} \begin{enumerate}[label={$S_{\arabic*}=$}] ')
for func in funcs:
	print('\item ')
	solution = set(sp.solve(eval(func), x))
	string =  '$'+   sp.latex(solution ) + '$; '
	print(string)
print(r'\end{enumerate} \end{multicols} ')
\end{pycode}
\end{proof}
\begin{proof}[solutions de l'exercice \ref{chapter07:exercice:equationE}]\hfill  
\begin{pycode}
from re import sub
import sympy as sp
from sympy.core.sympify import kernS

x, y, z, t , a, b ,A= sp.symbols('x y z t a b A', real=True)
k, m, n = sp.symbols('k m n', integer=True)
f, g, h = sp.symbols('f g h', cls=sp.Function)

# Create a list of functions to include in the table
funcs = [
'2*sp.pi*x-A',
'2*sp.pi*y/x-A',
'3*x-5-2*y',
'5*x+7*y-1',
'-3-2*x+9*y',
'12-2*y-5*x',
'5*y+3-2/x',
'7-3*y-5/x',
'7*y-2*x*y-5',
'2*y*x+3-3*x-7',
'5*y*x+3-7*x-2*y',
'2*y*x-3-3*x-5*y',
]
n=1
print(r'\begin{multicols}{4} \begin{enumerate}[label={\alph*)}] ')
for func in funcs:
	print('\item ')
	solution = set(sp.solve(eval(func), x))
	string =  '$'+   sp.latex(solution ) + '$; '
	print(string)
print(r'\end{enumerate} \end{multicols} ')
\end{pycode}
\end{proof}